\chapter{Nested Bethe Ansatz and Matrix Bethe--Yang Equations}
\label{ch:integrable-nested-bethe-ansatz-matrix-bethe-yang}

\ConventionsLorentzian

The algebraic Bethe ansatz of
Chapter~\ref{ch:integrable-algebraic-bethe-ansatz-transfer-matrices}
diagonalizes a transfer matrix by turning the Yang--Baxter equation into
commuting finite-dimensional operators.  This chapter applies the same
mechanism to non-diagonal finite-volume scattering.  The physical rapidities
are fixed by Bethe--Yang quantization, while auxiliary roots diagonalize the
internal transfer matrix.

\section{Matrix Bethe--Yang Quantization}

Let a massive integrable QFT have \(n\) stable particles with rapidities
\(\theta_1,\ldots,\theta_n\), masses \(m_{a_j}\), and internal species spaces
\(V_{a_j}\).  Write
\[
  \mathcal V_{\boldsymbol a}
  =
  V_{a_1}\otimes\cdots\otimes V_{a_n}.
\]
For a chosen particle \(j\), transporting it once around a spatial circle of
circumference \(L\) produces the phase
\(\exp(\ii m_{a_j}L\sinh\theta_j)\) and the ordered product of two-body
scatterings through the other particles.  In the \(R\)-matrix convention,
this ordered product is an internal transfer matrix
\[
  \mathcal T_j(\theta_j\mid\theta_1,\ldots,\widehat{\theta_j},\ldots,\theta_n)
  \in \operatorname{End}(\mathcal V_{\boldsymbol a}).
\]
The hat means that \(\theta_j\) is omitted from the list of spectator
rapidities.  The Yang--Baxter equation implies that the different
\(\mathcal T_j\)'s commute on the common domain where the scattering maps are
defined, after the usual care with identical particles and ordering
conventions.

\begin{definition}[Matrix Bethe--Yang state]
\label{def:matrix-bethe-yang-state}
A finite-volume \(n\)-particle matrix Bethe--Yang state is a vector
\(\Psi\in\mathcal V_{\boldsymbol a}\) and rapidities
\(\theta_1,\ldots,\theta_n\) such that \(\Psi\) is a simultaneous eigenvector
of all internal transfer matrices,
\[
  \mathcal T_j\Psi=\tau_j\Psi ,
\]
and
\begin{equation}
  \exp\!\left(\ii m_{a_j}L\sinh\theta_j\right)\tau_j=1,
  \qquad j=1,\ldots,n .
  \label{eq:matrix-bethe-yang}
\end{equation}
\end{definition}

When the scattering is diagonal, \(\tau_j\) is the product of scalar
two-particle phases and \eqref{eq:matrix-bethe-yang} reduces to the
Bethe--Yang equation in
Chapter~\ref{ch:integrable-thermodynamic-bethe-ansatz}.  In the non-diagonal
case the eigenvalue \(\tau_j\) is obtained by an auxiliary Bethe ansatz.

\section{Nested Roots for SU(N)}

Consider a homogeneous rational \(SU(N)\) chain with local space
\(V=\C^N\) in the fundamental representation and rational \(R\)-matrix
\[
  R(u)=u\,\id+\ii P .
\]
Choose a reference color \(1\).  A state with colors different from \(1\)
contains \(M_1\) first-level excitations.  If more than one non-reference
color is present, diagonalization of their internal scattering requires
further auxiliary roots.  For \(SU(N)\) there are \(N-1\) levels of roots
\[
  \{u_j^{(r)}\}_{j=1}^{M_r},\qquad r=1,\ldots,N-1 .
\]
The nested Bethe equations in the normalization compatible with
\eqref{eq:aba-xxx-bethe-equations} are
\begin{equation}
\begin{aligned}
  \left(
    \frac{u_j^{(1)}+\ii/2}{u_j^{(1)}-\ii/2}
  \right)^L
  &=
  \prod_{\substack{k=1\\ k\neq j}}^{M_1}
  \frac{u_j^{(1)}-u_k^{(1)}+\ii}
       {u_j^{(1)}-u_k^{(1)}-\ii}
  \prod_{k=1}^{M_2}
  \frac{u_j^{(1)}-u_k^{(2)}-\ii/2}
       {u_j^{(1)}-u_k^{(2)}+\ii/2},
\\
  1
  &=
  \prod_{\substack{k=1\\ k\neq j}}^{M_r}
  \frac{u_j^{(r)}-u_k^{(r)}+\ii}
       {u_j^{(r)}-u_k^{(r)}-\ii}
  \prod_{\sigma=\pm1}
  \prod_{k=1}^{M_{r+\sigma}}
  \frac{u_j^{(r)}-u_k^{(r+\sigma)}-\ii/2}
       {u_j^{(r)}-u_k^{(r+\sigma)}+\ii/2},
\\[-0.4em]
  &\hspace{19em} r=2,\ldots,N-2,
\\
  1
  &=
  \prod_{\substack{k=1\\ k\neq j}}^{M_{N-1}}
  \frac{u_j^{(N-1)}-u_k^{(N-1)}+\ii}
       {u_j^{(N-1)}-u_k^{(N-1)}-\ii}
  \prod_{k=1}^{M_{N-2}}
  \frac{u_j^{(N-1)}-u_k^{(N-2)}-\ii/2}
       {u_j^{(N-1)}-u_k^{(N-2)}+\ii/2}.
\end{aligned}
  \label{eq:su-n-nested-bethe}
\end{equation}
Products over an empty set are \(1\).  The energy of the fundamental rational
chain in the Hamiltonian normalization of
\eqref{eq:aba-xxx-energy-momentum} depends only on the first-level roots:
\[
  E=\sum_{j=1}^{M_1}\frac{1}{(u_j^{(1)})^2+1/4}.
\]

\begin{proposition}[Origin of the \(SU(N)\) nested equations]
\label{prop:su-n-nested-origin}
The equations \eqref{eq:su-n-nested-bethe} are the pole-cancellation
conditions for the nested transfer-matrix eigenvalue obtained by applying the
algebraic Bethe ansatz successively to the residual internal
\(SU(N-r)\)-color problem at levels \(r=1,\ldots,N-1\).
\end{proposition}

\begin{proof}
At the first level, the monodromy has one reference diagonal entry and an
\((N-1)\times(N-1)\) block acting on the non-reference colors.  Commuting the
reference diagonal entry through first-level creation operators produces the
first product in the first line of \eqref{eq:su-n-nested-bethe}.  The
unwanted terms are not scalars because the non-reference colors still scatter
among themselves.  Their coefficients are entries of an auxiliary transfer
matrix whose local space is \(\C^{N-1}\).  Diagonalizing this auxiliary
transfer matrix repeats the same construction with one fewer color.  At level
\(r\), roots at the same level repel with shift \(\pm\ii\), while roots at
adjacent levels couple with shift \(\pm\ii/2\).  Non-adjacent levels do not
share a simple root of \(A_{N-1}\), hence there is no factor.  Requiring every
apparent pole in the nested eigenvalue to have zero residue gives exactly
\eqref{eq:su-n-nested-bethe}.  The iteration terminates at \(r=N-1\), where
there is no further auxiliary color.
\end{proof}

The proof is finite-dimensional algebra.  Completeness of the Bethe vectors
for every chain length and every admissible representation is a separate
problem; it depends on boundary conditions, singular roots, and the treatment
of coincident solutions.

\section{Worked SU(3) Configuration}

For \(SU(3)\) there are first-level roots \(u_j\) and second-level roots
\(v_\alpha\).  The equations are
\begin{align}
  \left(\frac{u_j+\ii/2}{u_j-\ii/2}\right)^L
  &=
  \prod_{\substack{k=1\\ k\neq j}}^{M_1}
  \frac{u_j-u_k+\ii}{u_j-u_k-\ii}
  \prod_{\alpha=1}^{M_2}
  \frac{u_j-v_\alpha-\ii/2}{u_j-v_\alpha+\ii/2},
  \label{eq:su3-first-level}\\
  1
  &=
  \prod_{j=1}^{M_1}
  \frac{v_\alpha-u_j-\ii/2}{v_\alpha-u_j+\ii/2}.
  \label{eq:su3-second-level}
\end{align}
Set \(L=6\), \(M_1=2\), \(M_2=1\),
\[
  u_1=-\frac{\sqrt3}{2},\qquad
  u_2=\frac{\sqrt3}{2},\qquad
  v_1=0 .
\]
Let
\[
  z=\frac{u_2+\ii/2}{u_2-\ii/2}=\ee^{\ii\pi/3}.
\]
Then \(z^6=1\).  The second-level equation is
\[
  \frac{\sqrt3/2-\ii/2}{\sqrt3/2+\ii/2}
  \frac{-\sqrt3/2-\ii/2}{-\sqrt3/2+\ii/2}
  =
  z^{-1}z=1 .
\]
For the first root \(u_2\), the right side of
\eqref{eq:su3-first-level} is
\[
  \frac{2u_2+\ii}{2u_2-\ii}
  \frac{u_2-\ii/2}{u_2+\ii/2}
  =
  z\,z^{-1}=1,
\]
and the left side is \(z^6=1\).  The equation for \(u_1\) is the inverse of
the same equality.  The energy is
\[
  E=2\frac{1}{3/4+1/4}=2 .
\]
This example is deliberately small: every phase can be checked exactly, and
the second-level root is essential even though it carries no direct energy.

\section{Relativistic SU(N)-Symmetric Scattering}

For a relativistic integrable QFT whose particles transform in a fundamental
\(SU(N)\) multiplet, the two-body \(S\)-matrix often factorizes as
\[
  S(\theta)=S_0(\theta)\,\widehat R(\theta),
\]
where \(S_0\) is a scalar meromorphic factor and \(\widehat R\) is an
\(SU(N)\)-invariant rational or trigonometric matrix normalized to satisfy
unitarity and crossing in the physical rapidity variable \(\theta\).  The
finite-volume quantization separates into a scalar phase and a nested
transfer-matrix eigenvalue:
\begin{equation}
  mL\sinh\theta_j
  -\ii\sum_{k\neq j}\log S_0(\theta_j-\theta_k)
  -\ii\log \tau_j(\theta_1,\ldots,\theta_n;\{u^{(r)}\})
  =
  2\pi I_j .
  \label{eq:relativistic-nested-bethe-yang}
\end{equation}
The auxiliary roots \(\{u^{(r)}\}\) obey equations of the form
\eqref{eq:su-n-nested-bethe}, with rapidity-dependent shifts determined by
the matrix part \(\widehat R\).  The precise shift convention is part of the
chosen \(S\)-matrix normalization; changing it is a reparametrization of
auxiliary rapidities, not a change of physical content.

\begin{controlledapproximation}[Large-volume Bethe--Yang regime]
\label{ca:matrix-bethe-yang-large-volume}
Equation~\eqref{eq:relativistic-nested-bethe-yang} is a large-\(L\)
quantization condition for stable particles whose separations are much larger
than the interaction range set by the mass gap.  Corrections are
exponentially small in \(mL\) for a massive theory without massless channels,
provided no bound-state pole is pinched by the finite-volume contour.  Power
corrections and wrapping corrections require the TBA or mirror-channel
formalism of later chapters.
\end{controlledapproximation}

\section{Principal Chiral Model and Gross--Neveu Families}

The \(SU(N)\) principal chiral model has global
\(SU(N)_L\times SU(N)_R\) symmetry.  Its massive particles carry left and
right internal indices, so the matrix part of its scattering problem contains
two nested \(SU(N)\) systems, one for each chiral factor.  Gross--Neveu
families have a similar nested structure, with the details controlled by the
vector, spinor, or fundamental representation carried by the particles.  The
Bethe equations therefore contain physical rapidities together with auxiliary
roots attached to the Dynkin diagram of the relevant symmetry algebra.

The important conceptual point is that auxiliary roots are not additional
particles in the Hilbert space.  They parametrize the internal wavefunction of
a multiparticle state.  The energy and momentum are functions of the physical
rapidities; the auxiliary roots enter by selecting the transfer-matrix
eigenvalue that appears in the quantization condition.

\section{Status of Completeness and Continuum QFT}

For a finite chain, the algebraic Bethe ansatz constructs a large commuting
family of operators and many simultaneous eigenvectors.  Showing that these
vectors form a basis requires additional work.  For a continuum QFT the
logical burden is larger: one must also prove that the finite-volume spectrum
is approximated by the matrix Bethe--Yang equations in a controlled regime
and then control the thermodynamic or infinite-volume limit.  This monograph
therefore treats nested Bethe equations as exact algebraic data of the
factorized scattering framework and states explicitly whenever a spectral
completeness assertion is used as an additional theorem or assumption.
